\section{Second Domain: GloBI}
\label{app:globi}

We repeat the abstention test and the page edits on species interactions. This tests whether \clade{} depends on medicine, on MONDO, or on PrimeKG. It does not: the same pattern appears with a different graph, a different hierarchy and different text.

\subsection{Setup}

\paragraph{Facts and pages.}
Each fact links a consumer species to a partner species, such as a butterfly and the plant its larvae eat. The page is the consumer's Wikipedia article, cleaned as in \cref{app:data}. A page is \emph{named} if it gives the partner's Latin name, its abbreviated form (\emph{A.~plantaginifolia}) or a common name from Wikipedia redirects. It is \emph{family-only} if it gives only the partner's genus or family, and \emph{silent} if it gives neither. We measure \num{23694} edges over \num{5425} articles. Pages name the partner for only 3.9\% of edges and name only its genus or family for 7.1\% (\cref{M-tab:coverage}). Family-only pages are thus about twice as common as named ones.

\paragraph{Questions.}
We use one template per relation:
\begin{center}
\footnotesize
\begin{tabular}{@{}ll@{}}
\toprule
Relation & Question \\
\midrule
eats & Which species does X feed on? \\
preys on & Which species does X prey on? \\
parasite of & Which species is X a parasite of? \\
pollinates & Which plant does X pollinate? \\
\bottomrule
\end{tabular}
\end{center}
We take 75 questions per page type and relation, for 900 in all. Wrong options are other species of the partner's genus first, then of its family, and never another recorded partner of the consumer. For binomial answers, we fill any remaining slots with species from other genera, never with a bare genus name, so every option is a species. Most questions have four neighborhood options (833 of 900; \cref{tab:globi_sib}), and the main analysis keeps only these.

\begin{table}[t]
\centering
\footnotesize
\setlength{\tabcolsep}{3.5pt}
\begin{tabular}{@{}lrrrr@{}}
\toprule
Page type & eats & parasite of & pollinates & preys on \\
\midrule
Named & 93.3 & 94.7 & 94.7 & 89.3 \\
Family-only & 100.0 & 98.7 & 100.0 & 93.3 \\
Silent & 90.7 & 85.3 & 92.0 & 78.7 \\
\bottomrule
\end{tabular}
\caption{Share of GloBI questions (\%, 75 per cell) whose four wrong options are all from the partner's genus or family.}
\label{tab:globi_sib}
\end{table}

\subsection{Results per Model}

\paragraph{Outcomes.}
\Cref{tab:globi_models} gives outcomes for each model, closed book and with the page. Closed book, all three models decline most species questions; they know few of these facts. A named page lets them answer, and their answers are precise (84--95\%). A family-only page also lets them answer, but most of these answers are wrong: precision is 22--28\%, close to the 20\% chance rate of a five-way guess. A silent page makes them decline almost everything.

\begin{table}[t]
\centering
\footnotesize
\setlength{\tabcolsep}{2.4pt}
\begin{tabular}{@{}llrrrrrr@{}}
\toprule
 & & \multicolumn{2}{c}{Correct} & \multicolumn{2}{c}{Wrong} & \multicolumn{2}{c}{Declines} \\
\cmidrule(lr){3-4}\cmidrule(lr){5-6}\cmidrule(l){7-8}
Model & Page & cl. & page & cl. & page & cl. & page \\
\midrule
Qwen2.5-7B & named & 7.9 & 56.6 & 4.3 & 3.2 & 87.8 & 40.1 \\
 & family & 1.7 & 1.7 & 2.7 & 4.4 & 95.6 & 93.9 \\
 & silent & 3.5 & 0.4 & 1.2 & 1.2 & 95.4 & 98.5 \\
OLMo 2 7B & named & 20.8 & 69.9 & 24.7 & 12.9 & 54.5 & 17.2 \\
 & family & 10.9 & 12.6 & 22.8 & 34.4 & 66.3 & 53.1 \\
 & silent & 8.5 & 3.5 & 22.3 & 8.8 & 69.2 & 87.7 \\
Granite 3.1 8B & named & 12.2 & 60.2 & 16.8 & 7.2 & 71.0 & 32.6 \\
 & family & 9.2 & 3.4 & 14.6 & 12.2 & 76.2 & 84.4 \\
 & silent & 5.8 & 1.5 & 12.3 & 1.9 & 81.9 & 96.5 \\
\bottomrule
\end{tabular}
\caption{GloBI outcomes (\%) closed book (cl.) and with the page, on questions whose four wrong options all come from the neighborhood (279 named, 294 family-only, 260 silent).}
\label{tab:globi_models}
\end{table}

\paragraph{Page effects.}
\Cref{tab:globi_did} gives the page effect for each model. The effect is largest in OLMo~2, where a silent page raises declines by 18.5 points but a family-only page lowers them by 13.3, a difference of 31.7 points. Qwen2.5-7B declines almost every question on either page, so its effect is small but still significant. Granite's decline effect is not significant alone, but its wrong answers rise reliably. Pooled over the three models, $\Delta=-14.3$ points for declines and $+11.6$ for wrong answers, with only $+2.7$ for correct answers: in GloBI, the answers that the family name frees are mostly wrong.

\begin{table}[t]
\centering
\footnotesize
\setlength{\tabcolsep}{2.4pt}
\begin{tabular}{@{}llrlr@{}}
\toprule
Model & Outcome & $\Delta$ & [95\% CI] & $p$ \\
\midrule
Qwen2.5-7B & declines & $-4.8$ & [$-8.9$, $-0.8$] & .022 \\
 & wrong & $+1.7$ & [$-1.1$, $+4.8$] & .26 \\
OLMo 2 7B & declines & $-31.7$ & [$-40.6$, $-22.7$] & 2e-11 \\
 & wrong & $+25.0$ & [$+16.5$, $+33.4$] & 1e-8 \\
Granite 3.1 8B & declines & $-6.5$ & [$-14.0$, $+0.8$] & .097 \\
 & wrong & $+8.0$ & [$+1.5$, $+14.6$] & .019 \\
\midrule
Pooled & declines & $-14.3$ & [$-18.7$, $-9.6$] & 1e-9 \\
 & wrong & $+11.6$ & [$+7.5$, $+15.4$] & 3e-8 \\
\bottomrule
\end{tabular}
\caption{GloBI page effects (points), family-only minus silent (Eq.~\ref{M-eq:did}).}
\label{tab:globi_did}
\end{table}

\paragraph{Page edits and confidence.}
The per-model edit results are in \cref{tab:edit_models} and the confidence results in \cref{tab:conf}. In every model, the genus name in place of the answer draws fewer declines than deleting the sentence, and AUROC on family-only pages falls to 0.52--0.68.

\paragraph{Why GloBI is harder.}
Sibling species are harder to tell apart than sibling diseases. The models know few species facts closed book (\cref{tab:globi_models}), so when a page names only the genus, they guess inside the genus close to chance. This is why precision on family-only pages is near chance in GloBI (22--28\%) but closer to half in PrimeKG (45--59\%).
